\documentclass[11pt]{amsart}
\usepackage[T1]{fontenc}
\usepackage{lmodern,amsmath,amssymb,amsthm,mathtools,mathrsfs,microtype}
\usepackage[margin=1in]{geometry}
\usepackage{needspace}
\usepackage[hidelinks]{hyperref}
\numberwithin{equation}{section}
\newtheorem{theorem}{Theorem}[section]
\newtheorem{proposition}[theorem]{Proposition}
\newtheorem{lemma}[theorem]{Lemma}
\newcommand{\D}{\mathbb D}
\newcommand{\T}{\mathbb T}
\newcommand{\Hh}{\mathcal H}
\newcommand{\Id}{I}
\newcommand{\tr}{\operatorname{tr}}
\newcommand{\Dom}{\operatorname{Dom}}
\newcommand{\ran}{\operatorname{ran}}
\newcommand{\rank}{\operatorname{rank}}
\newcommand{\ip}[2]{\langle #1,#2\rangle}
\newcommand{\norm}[1]{\lVert #1\rVert}
\newcommand{\OO}{\mathcal O}
\newcommand{\VV}{\mathcal T}
\newcommand{\KK}{\mathcal K}
\title[The strict Dirichlet P\'olya inequality]{The strict Dirichlet P\'olya inequality on simply connected planar domains}
\date{}
\begin{document}
\begin{abstract}
We prove that the Dirichlet eigenvalues of every bounded simply connected
planar domain satisfy $|\Omega|\lambda_j(\Omega)>4\pi j$ for all $j\ge1$,
without any boundary regularity assumption. The proof associates with
conformal boundary curves a unitary path with positive trace-class
generator, injective at every positive radius, and integrated trace
equal to one half of the energy times the area. A boundary identity transfers Dirichlet
resonances to unitary eigenphases approaching a full turn. A Ky Fan
estimate controls any prescribed number of crossings by the corresponding
partial eigenvalue sum of the generator. The remaining spectral tail is
strictly positive and survives conformal exhaustion.
\end{abstract}
\maketitle

\section{Introduction}

Let $G\subset\mathbb R^2$ be a bounded nonempty open set.
The Dirichlet Laplacian $-\Delta_G^D$ is the self-adjoint operator in
$L^2(G)$ associated with the quadratic form
\[
 q_G[u]=\int_G|\nabla u|^2,\qquad \Dom q_G=H_0^1(G).
\]
We denote its eigenvalues, repeated according to multiplicity, by
$\lambda_1(G)\le\lambda_2(G)\le\cdots$, and write
\[
 N_G(E)=\#\{j:\lambda_j(G)\le E\}.
\]
We set $N_G(E)=0$ for $E\le0$. All counting functions in this paper
include the eigenvalues at the stated energy.

For a bounded planar domain with smooth boundary, Weyl's law
\cite{Weyl} gives
\[
 N_G(E)=\frac{|G|}{4\pi}E+o(E)\qquad(E\to\infty).
\]
P\'olya asked whether the leading term is an upper bound at every
energy, rather than only an asymptotic approximation. The conjecture
appears in his 1954 work on plausible reasoning \cite{Polya54}.
In its planar Dirichlet form, it asserts that
\[
 \lambda_j(G)\ge\frac{4\pi j}{|G|},\qquad j\ge1.
\]
The strict form replaces this inequality by $>$ at every index.

P\'olya proved the non-strict inequality for domains that tile the
plane \cite{Polya61}. The sharp inequalities of Berezin \cite{Berezin}
and Li--Yau \cite{LiYau} control spectral averages on general domains,
but do not give the conjectured bound for each eigenvalue. Freitas and
Salavessa \cite{FS} obtained families of non-tiling examples. Filonov,
Levitin, Polterovich, and Sher proved the conjecture for disks and
arbitrary circular sectors, and the Dirichlet conjecture for Euclidean
balls \cite{FLPSballs}. They subsequently treated planar annuli
\cite{FLPSannuli}. These results exploit special geometry; they do not
supply the inequality for an arbitrary simply connected planar domain.

Our main result gives the strict inequality for all bounded simply
connected planar domains, without a regularity assumption on the boundary.

\Needspace{12\baselineskip}
\begin{theorem}\label{thm:main}
Let $\Omega\subset\mathbb R^2$ be a bounded nonempty simply connected
domain. Its Dirichlet eigenvalues satisfy
\begin{equation}\label{eq:strict-eigenvalues}
 |\Omega|\lambda_j(\Omega)>4\pi j\qquad(j\ge1).
\end{equation}
Its counting function satisfies
\begin{equation}\label{eq:polya}
 N_\Omega(E)<\frac{|\Omega|E}{4\pi}\qquad(E>0).
\end{equation}
\end{theorem}

Here is the mechanism of the proof. At a fixed energy $k^2$, we
transport the unilateral shift along conformal boundary curves.
After correcting for the moving initial point and reversing the
resulting unitary operator, we obtain a path $U_r$ whose generator
$L_r=-iU_r^*\partial_rU_r$ is positive and trace class. Its integrated
trace is exactly
\[
 \int_0^1\tr L_r\,dr=\frac{k^2}{2}|\Omega|.
\]
An identity involving the Dirichlet-to-Neumann map relates a
Dirichlet eigenvalue to eigenvalues of $U_r$ approaching $1$ through
the lower half of the unit circle. Counting crossings of $e^{i\alpha}$
for $\alpha<2\pi$ and then letting $\alpha\uparrow2\pi$ recovers
the coefficient $4\pi$.

Strictness requires more than this full-trace estimate. If at least
$m$ crossings occur, we bound their contribution using only the sum
of the $m$ largest eigenvalues of $L_r$. The difference between the
full trace and this partial sum is positive, since $L_r$ is injective
on an infinite-dimensional space for every $r>0$. We retain that
difference through the limit in $\alpha$ and through conformal
exhaustion. This gives the strict inequality even when the boundary
is not regular.

The distinction between an eigenvalue approaching $1$ and an actual
eigenvalue at $1$ is essential. A related distinction occurs in the
inside--outside spectral duality of Eckmann and Pillet \cite{EP}.
We do not identify the boundary transport used here with their
scattering matrix. Instead, we prove the required boundary identity
directly on traces of Herglotz waves. This also specifies the domain
of the unbounded Cayley transform. The density argument is classical
in scattering theory; see Colton and Kress \cite{CK}. We retain its
short proof in the precise trace topology needed below.

Section~\ref{sec:phase} gives the eigenphase estimate with a prescribed
number of crossings. Section~\ref{sec:holonomy} constructs the boundary
transport and proves the trace identity and positivity of the spectral
tail. Section~\ref{sec:boundary} establishes the boundary identity on
Herglotz traces. Section~\ref{sec:poles} relates Dirichlet eigenvalues
to eigenphases near $2\pi$. Section~\ref{sec:completion} completes the
counting argument and passes the positive remainder to arbitrary
simply connected domains.

Throughout, inner products are conjugate-linear in the first variable.
We identify $\mathbb R^2$ with $\mathbb C$ and put
$\D=\{z:|z|<1\}$ and $\T=\mathbb R/(2\pi\mathbb Z)$.
Boundary spaces on $\T$ use $d\theta$; the direction space of
Herglotz waves uses $d\phi/(2\pi)$. The pairing
$\ip h{\Lambda h}$ is the $H^{1/2}$--$H^{-1/2}$ duality pairing,
consistent with the $L^2(d\theta)$ inner product for smooth data.
If $h_n$ are the coefficients in the orthonormal basis
$(2\pi)^{-1/2}e^{in\theta}$, our Sobolev norm is
$\norm h_{H^s(\T)}^2=\sum_{n\in\mathbb Z}(1+n^2)^s|h_n|^2$.
We write $\mathcal B(\Hh)$ for the bounded operators on $\Hh$,
and $\Id$ for the identity. The symbols $\sigma$, $\ker$, $\ran$,
$\rank$, and $\Dom$ denote spectrum, kernel, range, rank, and domain.
The symbols $\mathcal S_1$ and $\mathcal S_2$ denote the trace-class
and Hilbert--Schmidt ideals, respectively, and $\tr$ denotes the trace.
For bounded operators, $[A,B]=AB-BA$.

\section{Eigenphase estimates for positive unitary paths}\label{sec:phase}

For a positive trace-class operator $A$, let
$\mu_1(A)\ge\mu_2(A)\ge\cdots\ge0$ be its eigenvalues, repeated
according to multiplicity and supplemented by zeros when necessary.
For every integer $m\ge1$, set
\begin{equation}\label{eq:fan-def}
 \Phi_m(A)=\sum_{\ell=1}^m\mu_\ell(A),\qquad
 \tau_m(A)=\tr A-\Phi_m(A).
\end{equation}
We use Ky Fan's variational principle \cite{Fan} in the form
\begin{equation}\label{eq:fan-variational}
 \Phi_m(A)=\sup\{\tr(PA):P=P^*=P^2,\ \rank P\le m\}.
\end{equation}
In particular, for positive trace-class $A$ and $B$,
\begin{equation}\label{eq:fan-continuity}
 |\Phi_m(A)-\Phi_m(B)|\le\norm{A-B}_{\mathcal S_1},\qquad
 |\tau_m(A)-\tau_m(B)|\le2\norm{A-B}_{\mathcal S_1}.
\end{equation}
Indeed, the first bound follows by taking the supremum in
\eqref{eq:fan-variational} and using
$|\tr(P(A-B))|\le\norm{A-B}_{\mathcal S_1}$; the second also uses
continuity of the trace. If $A$ is injective on an infinite-dimensional
space, then $\tau_m(A)>0$ for every finite $m$.

For $0<\alpha<2\pi$, a crossing of $e^{i\alpha}$ is a time at which
it is an eigenvalue of the unitary path. Its multiplicity is the
dimension of that eigenspace. All such crossings in the next lemma
are isolated and have increasing local eigenphases.

\begin{lemma}\label{lem:phase}
Let $U:[0,b]\to\mathcal B(\Hh)$ be a continuously differentiable
unitary path on a separable Hilbert space, with $U(0)=\Id$.
Suppose
\[
 L(t)=-iU(t)^*U'(t)\ge0
\]
is trace class and trace-norm continuous. Assume that $U$ is locally
real analytic in operator norm for $t>0$, and that
$\ip v{L(t)v}>0$ for every $v\ne0$ and $t>0$.
Set $B=\int_0^b\tr L(t)\,dt$.
Then $U(t)-\Id$ is trace class, and the following statements hold.

\smallskip\noindent
\textup{(i)} For $0<\alpha<2\pi$, the number $c_\alpha$ of crossings
of $e^{i\alpha}$ in $(0,b)$, counted with multiplicity, is finite,
and all crossings are positive. For every integer $m$ with
$1\le m\le c_\alpha$,
\begin{equation}\label{eq:fan-cost}
 \alpha m\le\int_0^b\Phi_m(L(t))\,dt.
\end{equation}
In particular,
\begin{equation}\label{eq:cost}
 \alpha c_\alpha\le B.
\end{equation}

\smallskip\noindent
\textup{(ii)} At each fixed time, $U(t)$ has only finitely many
eigenvalues in $\{e^{i\theta}:\pi<\theta<2\pi\}$, with multiplicity.

\smallskip\noindent
\textup{(iii)} Let $0<a<b_1\le b$ and $\pi<\alpha<2\pi$.
Suppose $U(a)$ has no spectrum with argument in $[\alpha,2\pi)$,
$e^{i\alpha}\notin\sigma(U(b_1))$, and $U(b_1)$ has at least $m$
eigenvalues with arguments in $(\alpha,2\pi)$. There are at least
$m$ crossings of $e^{i\alpha}$ in $(a,b_1)$.
\end{lemma}

\begin{proof}
Since $U'=iUL$, integration in trace norm gives
$U(t)-\Id\in\mathcal S_1$. Every spectral point other than $1$
is therefore an isolated eigenvalue of finite multiplicity.
Near such a point, analytic Riesz projections and a local logarithm
reduce the spectral problem to an analytic Hermitian matrix family;
we use the standard perturbation results in
\cite[Chapters II and VII]{Kato}. For a normalized eigenvector
branch with eigenvalue $e^{i\theta(t)}$,
\[
 \theta'(t)=\ip{v(t)}{L(t)v(t)}.
\]
At a multiple eigenvalue, the eigenphase derivatives are the
eigenvalues of the compression of $L(t)$ to the eigenspace.
This compression is positive definite. Hence every crossing away
from $1$ is locally isolated and positive.

Choose finite-rank orthogonal projections $P_n\uparrow\Id$ strongly,
put $L_n=P_nLP_n$, and let $U_n$ be the unitary solution on
$P_n\Hh$ of $U_n'=iU_nL_n$, $U_n(0)=\Id$; extend $U_n$ by the
identity on $P_n^\perp\Hh$. The image of $[0,b]$ under $L$ is
compact in $\mathcal S_1$, so
\[
 \sup_t\norm{L_n(t)-L(t)}_{\mathcal S_1}\longrightarrow0.
\]
Moreover,
\[
 \frac{d}{dt}(U_nU^*)=iU_n(L_n-L)U^*.
\]
Integration in trace norm therefore gives
\begin{equation}\label{eq:approx}
 \sup_t\norm{U_n(t)-U(t)}_{\mathcal S_1}
 \le\int_0^b\norm{L_n(t)-L(t)}_{\mathcal S_1}\,dt
 \longrightarrow0.
\end{equation}
The differential equations also give uniform operator-norm
convergence of the first derivatives.

In finite dimension, the eigenphases admit continuous,
absolutely continuous, nondecreasing lifts $\theta_{n,j}$ starting
at zero. They may be obtained from local Hermitian logarithms and
continued along the interval. At almost every time their derivatives
are diagonal entries of $L_n(t)$ in an orthonormal eigenbasis; at a
multiple eigenvalue one diagonalizes its compression. In particular,
\[
 \sum_j\theta_{n,j}(b)=\int_0^b\tr L_n(t)\,dt\le B.
\]
Fix any $m$ crossings of $e^{i\alpha}$, counted with multiplicity.
Let $q_j$ be the number of selected crossings belonging to the
$j$th lifted phase, and put $J=\{j:q_j>0\}$. Then
$\sum_{j\in J}q_j=m$ and $|J|\le m$. Monotonicity gives
\[
 \theta_{n,j}(b)\ge\alpha+2\pi(q_j-1)\ge\alpha q_j
 \qquad(j\in J).
\]
This remains true when the selected crossings omit earlier crossings
of the same branch. At almost every time, the eigenvectors belonging
to the branches in $J$ form an orthonormal family. The variational
principle \eqref{eq:fan-variational} therefore yields
\begin{equation}\label{eq:finite-fan-cost}
 \alpha m\le\sum_{j\in J}\theta_{n,j}(b)
 =\int_0^b\sum_{j\in J}\theta_{n,j}'(t)\,dt
 \le\int_0^b\Phi_m(L_n(t))\,dt\le B.
\end{equation}
Taking all crossings proves the finite-dimensional bound $B/\alpha$.
More generally, \eqref{eq:finite-fan-cost} applies whenever at least
$m$ crossings are available; they need not belong to distinct branches.

We explain the passage to $U$ without following eigenvalues through
$1$. Fix a crossing time $t_0\in(0,b)$ of multiplicity $q$.
Choose a small arc about $e^{i\alpha}$, disjoint from $1$, whose
spectral projection at $t_0$ has rank $q$. On a sufficiently small
closed time interval about $t_0$, that projection keeps rank $q$.
After identifying its ranges with a fixed space, the local logarithm
relative to $e^{i\alpha}$ is negative definite at the left endpoint
and positive definite at the right endpoint. By \eqref{eq:approx},
the corresponding projections for $U_n$ have the same rank, and
the endpoint logarithms have the same inertias for large $n$.
Their eigenvalues must therefore cross zero with total signed
multiplicity $q$. The positivity of $L_n$ on $P_n\Hh$ for positive
times makes these crossings positive. The same construction applies
to any finite collection of crossing times, using disjoint time
intervals. Their total multiplicity is at most $B/\alpha$.
This proves both finiteness and \eqref{eq:cost}. If $m\le c_\alpha$,
choose finitely many crossing windows containing at least $m$
crossings of $U$, with multiplicity. For all sufficiently large $n$
the path $U_n$ has at least $m$ crossings in those windows.
Apply \eqref{eq:finite-fan-cost} to $m$ of them and use
\eqref{eq:fan-continuity} and the uniform trace-norm convergence
of $L_n$ to $L$. This proves \eqref{eq:fan-cost}. Both estimates
remain valid for interior crossings when an endpoint of the full
path is itself a crossing.

At a fixed time, any finite collection of eigenvalues of $U(t)$
in the open lower semicircle persists, with multiplicity, under
\eqref{eq:approx}. Each corresponding eigenvalue of $U_n(t)$ has
a nonnegative lifted phase greater than $\pi$. The collection
therefore has size at most $B/\pi$. Since the collection was
arbitrary, this proves (ii).

For (iii), choose $\beta<2\pi$ above the arguments of $m$ selected
eigenvalues at $b_1$, with $e^{i\beta}$ outside both endpoint spectra.
The spectral projection onto the arc $(\alpha,\beta)$ has finite
rank. Between crossings of the endpoints its rank is constant.
Part (i) gives only finitely many endpoint crossings, and their
positive directions give the rank balance
\[
 n_{(\alpha,\beta)}(b_1)-n_{(\alpha,\beta)}(a)
 =c_\alpha(a,b_1)-c_\beta(a,b_1).
\]
Here $n$ denotes spectral multiplicity in the indicated arc, and
$c$ denotes crossing multiplicity in the indicated time interval.
The initial rank is zero, the final rank is at least $m$, and
$c_\beta(a,b_1)\ge0$. This proves (iii).
\end{proof}

\section{Boundary transport and the area identity}\label{sec:holonomy}

Until Section~\ref{sec:completion}, assume that
$\Omega=F(\D)$, where $F$ is holomorphic and injective on a
neighborhood of $\overline\D$. Fix $k>0$, and define
\[
 \Omega_r=F(r\D),\qquad \gamma_r(\theta)=F(re^{i\theta}),
 \qquad 0<r\le1.
\]
The curves are positively oriented. We write
$\gamma_r'=\partial_\theta\gamma_r$; radial derivatives of
operators will be written as $\partial_r$.

Let $\Hh=\ell^2(\mathbb N_0)$, where $\mathbb N_0=\{0,1,\ldots\}$, with standard basis $(e_n)_{n\ge0}$,
and let $Se_n=e_{n+1}$. If $P_0v=e_0\ip{e_0}v$, then
$[S^*,S]=P_0$. Define the unitary transport $W_r(\theta)$ by
\begin{equation}\label{eq:W}
 \partial_\theta W_r
 =-\frac{ik}{2}(\gamma_r'S^*+\overline{\gamma_r'}S)W_r,
 \qquad W_r(0)=\Id,
 \qquad V_r=W_r(2\pi).
\end{equation}
Writing $\gamma_r=x+iy$, set
\[
 X=-\frac{ik}{2}(S+S^*),\qquad Y=\frac{k}{2}(S^*-S).
\]
These are bounded skew-adjoint operators, and
$[X,Y]=ik^2P_0/2$. Put
\[
 C_\theta=x_\theta X+y_\theta Y,\qquad
 C_r=x_rX+y_rY,\qquad B_r=C_r(r,0)=C_r(r,2\pi).
\]
The moving initial point is removed by solving
$\partial_rR_r=B_rR_r$, $R_0=\Id$, and setting
\begin{equation}\label{eq:based}
 H_r=R_r^*V_rR_r,\qquad U_r=H_r^*.
\end{equation}
The loop at $r=0$ is constant, so $U_0=\Id$.

\begin{proposition}\label{prop:area}
The path $U_r$ satisfies the assumptions of Lemma~\ref{lem:phase}.
Its generator $L_r=-iU_r^*\partial_rU_r$ satisfies, for $0<b\le1$,
\begin{equation}\label{eq:area}
 \int_0^b\tr L_r\,dr=\frac{k^2}{2}|\Omega_b|.
\end{equation}
For every integer $m\ge1$ and every $r>0$,
\begin{equation}\label{eq:positive-tail}
 \tau_m(L_r)>0.
\end{equation}
Writing $L_{r,k}$ when the energy parameter varies, the map
$(r,k)\mapsto L_{r,k}$ is jointly continuous in trace norm for
$r\in[0,1]$ and $k>0$.
\end{proposition}

\begin{proof}
Differentiating \eqref{eq:W} yields
\[
 \partial_\theta(W_r^*\partial_rW_r-W_r^*C_rW_r)
 =W_r^*(\partial_rC_\theta-\partial_\theta C_r+[C_\theta,C_r])W_r.
\]
The mixed derivatives cancel. Since
\[
 J=x_ry_\theta-x_\theta y_r=r|F'(re^{i\theta})|^2,
 \qquad [C_\theta,C_r]=-J[X,Y],
\]
integration around the boundary gives
\[
 V_r^*\partial_rV_r
 =V_r^*B_rV_r-B_r-\frac{ik^2}{2}
       \int_0^{2\pi}J W_r^*P_0W_r\,d\theta.
\]
Differentiating \eqref{eq:based} cancels the first two terms:
\begin{equation}\label{eq:Q}
 H_r^*\partial_rH_r=-iQ_r,\qquad
 Q_r=\frac{k^2}{2}R_r^*
       \left(\int_0^{2\pi}J W_r^*P_0W_r\,d\theta\right)R_r.
\end{equation}
Thus $L_r=H_rQ_rH_r^*\ge0$. The integral in \eqref{eq:Q} is a
trace-norm integral of positive rank-one operators.
Taking the trace in \eqref{eq:Q} gives
\[
 \tr L_r=\frac{k^2}{2}\int_0^{2\pi}
       r|F'(re^{i\theta})|^2\,d\theta.
\]
Integration from $0$ to $b$ proves \eqref{eq:area} by the conformal area formula. The same representation gives
trace-norm continuity, including at $r=0$.

Suppose $r>0$ and $\ip v{Q_rv}=0$.
The zeroth component of $f(\theta)=W_r(\theta)R_rv$ vanishes for
every $\theta$, since $J>0$. Its differential equation is
$0=f_0'=-(ik/2)\gamma_r'f_1$. Since $\gamma_r'\ne0$, we obtain
$f_1=0$. Successively applying the component equations gives
$f_n=0$ for every $n$, and hence $v=0$. Unitary conjugation proves
the required strict positivity of $L_r$. In particular, $L_r$
is injective on the infinite-dimensional space $\Hh$. Its compactness
then implies \eqref{eq:positive-tail} by \eqref{eq:fan-def}.

The coefficients of both operator differential equations are locally
real analytic in $r$, uniformly in $\theta$. The uniformly convergent
Picard series on compact parameter neighborhoods gives operator-norm
real analyticity of $W_r,V_r,R_r,U_r$. At $r=0$ the same equations
and \eqref{eq:Q} give continuous differentiability. This verifies
the remaining assumptions of Lemma~\ref{lem:phase}.

Continuous dependence in the same differential equations gives joint
operator-norm continuity of $W_r(\theta)$, $R_r$, and $H_r$ in
$(r,k)$, uniformly in $\theta$ on compact parameter sets.
The integrand in \eqref{eq:Q} is a continuous rank-one operator
in trace norm. That formula and $L_r=H_rQ_rH_r^*$ give the stated
joint trace-norm continuity, including at $r=0$ where $J=0$.
\end{proof}

\section{Herglotz waves and the boundary identity}\label{sec:boundary}

\subsection{Density of boundary traces}
For $a\in L^2(\T,d\phi/(2\pi))$, define the Herglotz wave
\begin{equation}\label{eq:herglotz}
 u_a(z)=\int_0^{2\pi}a(\phi)
 \exp\left[-\frac{ik}{2}(ze^{-i\phi}+\bar z e^{i\phi})\right]
 \frac{d\phi}{2\pi}.
\end{equation}
It is smooth on $\mathbb R^2$ and satisfies $(\Delta+k^2)u_a=0$.
We use the following standard consequence of uniqueness for the
exterior Helmholtz problem; compare \cite{CK,EP}.

\begin{lemma}\label{lem:density}
Let $G$ be a bounded smooth Jordan domain. If
$k^2\notin\sigma(-\Delta_G^D)$, the traces of the waves
\eqref{eq:herglotz} are dense in $H^{1/2}(\partial G)$.
\end{lemma}

\begin{proof}
Let a continuous linear functional represented by
$g\in H^{-1/2}(\partial G)$ annihilate all the traces in question.
Form the outgoing single-layer potential with density $g$ and
fundamental solution $(i/4)H_0^{(1)}(k|x-y|)$, where $H_0^{(1)}$
is the Hankel function of the first kind and order zero.
The far field is, up to a nonzero constant, the pairing of $g$
with plane waves. The annihilation assumption for arbitrary
$L^2$ direction densities makes this far field zero in $L^2$
and then in every direction by continuity.

Rellich's uniqueness theorem, or the outgoing Fourier--Hankel
expansion outside a large disk, makes the potential zero there.
Unique continuation and connectedness of the exterior make it zero
throughout $\mathbb R^2\setminus\overline G$.
The single-layer mapping
$H^{-1/2}(\partial G)\to H^1_{\rm loc}(\mathbb R^2)$ ensures that
the interior and exterior Dirichlet traces agree; see
\cite[Section 3]{EP}. Thus the interior solution has zero boundary
trace and is zero by the nonresonance assumption. Applying
$-\Delta-k^2$ distributionally now gives
$g\delta_{\partial G}=0$, and hence $g=0$.
The Hahn--Banach theorem proves the assertion.
\end{proof}

Smooth boundary parametrizations preserve $H^{1/2}$, so the lemma
also applies to the traces pulled back to $\T$.
We call a radius $r$ nonresonant when
$k^2\notin\sigma(-\Delta_{\Omega_r}^D)$.

\subsection{The Cayley transform}

\begin{lemma}\label{lem:injective}
At every nonresonant radius,
$\ker(V_r-\Id)=\{0\}$.
\end{lemma}

\begin{proof}
Suppress $r$. Embed $\Hh$ as the nonnegative Fourier modes in
$\ell^2(\mathbb Z)=L^2(\T,d\phi/(2\pi))$, and let $Z$ be
multiplication by $\zeta=e^{i\phi}$. The full-space transport is
\[
 M(\theta)=\exp\left[-\frac{ik}{2}
 \big((\gamma(\theta)-\gamma(0))Z^*
 + (\overline{\gamma(\theta)}-\overline{\gamma(0)})Z\big)\right].
\]
Its generators commute, and $M(0)=M(2\pi)=\Id$.
Extend $f(\theta)=W(\theta)v$ by zero to the negative indices.
Then
\[
 f'=-\frac{ik}{2}(\gamma'Z^*+\overline{\gamma'}Z)f
       +\frac{ik}{2}\gamma'f_0e_{-1}.
\]
Variation of constants in the bilateral space gives
\[
 (V-\Id)v=\frac{ik}{2}\int_0^{2\pi}
       M(\theta)^*e_{-1}\gamma'(\theta)f_0(\theta)\,d\theta.
\]
If $Vv=v$, multiplication by $\zeta$ and by a nonvanishing
direction-dependent phase yields
\[
 \int_0^{2\pi}\gamma'(\theta)f_0(\theta)
       e^{ik\operatorname{Re}(\gamma(\theta)e^{-i\phi})}\,d\theta=0.
\]
This equality holds for every direction by continuity.
Changing a direction to its opposite shows that the continuous
linear functional
$h\mapsto\int\gamma'f_0h\,d\theta$ annihilates all Herglotz traces.
Lemma~\ref{lem:density} gives $\gamma'f_0=0$. The component equations
in \eqref{eq:W} then give $f_n=0$ successively, so $v=0$.
\end{proof}

At a nonresonant radius define
\begin{equation}\label{eq:cayley}
 \KK_r=i(\Id+V_r)(\Id-V_r)^{-1},
 \qquad \Dom\KK_r=\ran(\Id-V_r).
\end{equation}
Normality of $V_r$ and Lemma~\ref{lem:injective} make this domain
dense. The usual Cayley-transform construction gives a self-adjoint
operator. Directly, its two nonreal resolvents are
\[
 (\KK_r+i)^{-1}=\frac{\Id-V_r}{2i},\qquad
 (\KK_r-i)^{-1}=\frac{(\Id-V_r)V_r^*}{2i}.
\]
They are compact by Lemma~\ref{lem:phase} and
Proposition~\ref{prop:area}, so $\KK_r$ has compact resolvent.
By \eqref{eq:based}, an eigenvalue $e^{i\theta}$ of $U_r$, with
$0<\theta<2\pi$, corresponds, with the same multiplicity, to
\begin{equation}\label{eq:cot}
 \cot(\theta/2)\in\sigma(\KK_r).
\end{equation}
Lemma~\ref{lem:phase}(ii) implies that $\KK_r$ has only finitely
many negative eigenvalues. At each fixed nonresonant radius there
is therefore a gap in the spectrum of $U_r$ immediately below
argument $2\pi$. No uniform gap in $r$ is being asserted.

\subsection{A Dirichlet-to-Neumann identity}

Define $\OO_r:\Hh\to L^2(\T,d\theta)$ by
\[
 (\OO_rv)(\theta)=\gamma_r'(\theta)\ip{e_0}{W_r(\theta)v}.
\]
Its adjoint and Hilbert--Schmidt norm are
\begin{equation}\label{eq:Ostar}
 \OO_r^*h=\int_0^{2\pi}W_r(s)^*e_0\overline{\gamma_r'(s)}h(s)\,ds,
 \qquad \norm{\OO_r}_{\mathcal S_2}^2
       =\int_0^{2\pi}|\gamma_r'|^2\,d\theta.
\end{equation}
Also define the Volterra integral operator
\[
 (\VV_rh)(\theta)=\int_0^\theta
 \gamma_r'(\theta)\overline{\gamma_r'(s)}
 \ip{e_0}{W_r(\theta)W_r(s)^*e_0}h(s)\,ds.
\]
Comparing the kernels on the two triangles gives
\begin{equation}\label{eq:Volterra}
 \VV_r+\VV_r^*=\OO_r\OO_r^*.
\end{equation}
These operators are bounded locally uniformly in $r>0$.
The cut at $\theta=0$ causes no difficulty for these $L^2$ operators.

At a nonresonant radius, let
$\Lambda_r:H^{1/2}(\T)\to H^{-1/2}(\T)$ be the
Dirichlet-to-Neumann map with conormal density:
\[
 (\Lambda_rh)(\theta)=|\gamma_r'(\theta)|
       \partial_\nu u(\gamma_r(\theta)),\qquad
 (\Delta+k^2)u=0,\qquad u\circ\gamma_r=h,
\]
where $\nu$ is the outward unit normal.
Conformal pullback identifies $\Lambda_r$ with the disk
Dirichlet-to-Neumann map for
\begin{equation}\label{eq:Ar}
 \mathscr A_r=-\Delta-k^2w_r,\qquad
 w_r(z)=r^2|F'(rz)|^2.
\end{equation}
Its Dirichlet realization $A_r$ has domain
$H^2(\D)\cap H_0^1(\D)$. We distinguish the differential expression
$\mathscr A_r$ from the operator $A_r$ on that domain.

Put
\[
 \mathfrak r_r[h]=\ip h{i\partial_\theta h}
       +\frac{ik^2}{4}\ip h{(\VV_r-\VV_r^*)h}.
\]
This is a Hermitian quadratic form on $H^{1/2}(\T)$.

\begin{proposition}\label{prop:identity}
At a nonresonant radius, every Herglotz trace $h=u_a\circ\gamma_r$
satisfies $\OO_r^*h\in\Dom\KK_r$ and
\begin{equation}\label{eq:identity}
 \Lambda_rh=i\partial_\theta h
       +\frac{ik^2}{4}(\VV_r-\VV_r^*)h
       +\frac{k^2}{4}\OO_r\KK_r\OO_r^*h.
\end{equation}
For every compact interval $I\subset(0,1]$, there is a constant
$C_I$ such that
\begin{equation}\label{eq:remainder}
 |\mathfrak r_r[h]|\le C_I\norm h_{H^{1/2}(\T)}^2
       \qquad(r\in I,\ h\in H^{1/2}(\T)).
\end{equation}
\end{proposition}

\begin{proof}
Write $\partial=(\partial_x-i\partial_y)/2$ and
$\bar\partial=(\partial_x+i\partial_y)/2$, and set
\[
 f_n(z)=(2i/k)^{n+1}\partial^{n+1}u_a(z),\qquad n\ge0.
\]
These are the Fourier coefficients of indices $n+1$ of
$a(\phi)e^{-ik\operatorname{Re}(ze^{-i\phi})}$.
Parseval's identity gives $f(z)\in\Hh$ and
$\norm{f(z)}\le\norm a_{L^2(d\phi/(2\pi))}$.
Differentiation is valid in $\Hh$, because the differentiated
multipliers are bounded on the direction space. We have
\[
 \partial f_n=-\frac{ik}{2}f_{n+1},\qquad
 \bar\partial f_n=-\frac{ik}{2}f_{n-1},\qquad f_{-1}=u_a.
\]
On the boundary this gives
\[
 \partial_\theta f=-\frac{ik}{2}
       (\gamma_r'S^*+\overline{\gamma_r'}S)f
       -\frac{ik}{2}\overline{\gamma_r'}h e_0.
\]
If $v=f(\gamma_r(0))$, variation of constants gives
\[
 f(\theta)=W_r(\theta)v-\frac{ik}{2}W_r(\theta)
       \int_0^\theta W_r(s)^*e_0\overline{\gamma_r'(s)}h(s)\,ds.
\]
Periodicity and \eqref{eq:Ostar} imply
\begin{equation}\label{eq:endpoint}
 (\Id-V_r)v=-\frac{ik}{2}V_r\OO_r^*h.
\end{equation}
In particular,
\[
 \OO_r^*h=\frac{2i}{k}(V_r^*-\Id)v
       =\frac{2i}{k}(\Id-V_r)V_r^*v\in\Dom\KK_r.
\]
This establishes the domain assertion in \eqref{eq:cayley} before an inverse is used.

For a positively oriented curve, the outward normal is
$-i\gamma_r'/|\gamma_r'|$. A direct calculation therefore gives
\[
 |\gamma_r'|\partial_\nu u_a=i\partial_\theta h-k\gamma_r'f_0,
 \qquad \gamma_r'f_0=\OO_rv-\frac{ik}{2}\VV_rh.
\]
Nonresonance identifies $u_a$ with the solution defining $\Lambda_r$.
Using \eqref{eq:endpoint}, we obtain
\[
 \Lambda_rh=i\partial_\theta h
       +\frac{ik^2}{2}\OO_r(\Id-V_r)^{-1}V_r\OO_r^*h
       +\frac{ik^2}{2}\VV_rh.
\]
All inverses act on the ranges just established. On this domain,
\[
 i(\Id-V_r)^{-1}V_r=\tfrac12\KK_r-\tfrac i2\Id.
\]
Equation \eqref{eq:Volterra} proves \eqref{eq:identity}.

For the orthonormal Fourier basis $(2\pi)^{-1/2}e^{in\theta}$,
\[
 |\ip h{i\partial_\theta h}|
 \le\sum_{n\in\mathbb Z}|n||h_n|^2
 \le\norm h_{H^{1/2}(\T)}^2.
\]
The operator $i(\VV_r-\VV_r^*)$ is bounded and self-adjoint on
$L^2$, with locally uniform bounds given by its kernel.
This proves \eqref{eq:remainder}, including at resonant radii,
since no Dirichlet resolvent occurs in $\mathfrak r_r$.
\end{proof}

The last term in \eqref{eq:identity} is asserted only on Herglotz
traces. Density alone would not justify extending a composition
with the unbounded operator $\KK_r$ to all of $H^{1/2}(\T)$.

\section{Dirichlet eigenvalues and unitary eigenphases}\label{sec:poles}

We first record the sign and multiplicity of the singularity of
$\Lambda_r$ as a Dirichlet eigenvalue is approached from smaller
domains. The proof uses finite-dimensional spectral projections
rather than a choice of varying eigenfunctions.

\begin{lemma}\label{lem:pole}
Suppose $k^2$ is a Dirichlet eigenvalue of $\Omega_{r_*}$ of
multiplicity $m$, where $r_*>0$. There are linearly independent
$h_1,\ldots,h_m\in C^\infty(\T)$ and numbers $M_r\to+\infty$ as
$r\uparrow r_*$ such that, for every sufficiently close nonresonant
$r<r_*$,
\begin{equation}\label{eq:pole}
 \ip{h_c}{\Lambda_rh_c}\le-M_r|c|^2,
 \qquad h_c=\sum_{j=1}^m c_jh_j,\qquad c\in\mathbb C^m.
\end{equation}
\end{lemma}

\begin{proof}
Every ordered eigenvalue $\lambda_j(\Omega_r)$ is strictly decreasing
in $r$. Indeed, suppose $r<s$ and the two $j$th eigenvalues equal
$\lambda$. The first $j$ eigenfunctions on $\Omega_r$, extended by
zero, span a trial space on $\Omega_s$ with Rayleigh quotients at
most $\lambda$. It contains a nonzero vector orthogonal to the
spectral subspace of $\Omega_s$ below $\lambda$, whose dimension is
at most $j-1$. The spectral theorem forces that vector to be an
eigenfunction on $\Omega_s$ at $\lambda$. It vanishes on the
nonempty open set $\Omega_s\setminus\overline{\Omega_r}$,
contrary to unique continuation.

These eigenvalues are also continuous for $r>0$: use the min--max
principle on the fixed disk with denominator
$\int_\D w_r|u|^2$. On a compact positive-radius interval the
weights converge uniformly and are bounded away from zero.
By conformal invariance of the Dirichlet energy, the number of
negative eigenvalues of $A_r$ in \eqref{eq:Ar} is
$\#\{j:\lambda_j(\Omega_r)<k^2\}$, and its nullity is the
multiplicity of $k^2$.

The bounded potentials in $A_r$ vary continuously in $L^\infty$,
so the operators vary in norm resolvent. Let $P_r$ be the Riesz
projection of the spectral cluster near zero that equals
$\ker A_{r_*}$ at $r_*$. Its rank is $m$, and $P_r\to P_*$ in
operator norm. Strict monotonicity and continuity of the domain
eigenvalues show that the negative index below $r_*$ equals the
negative index at $r_*$. The spectrum outside the cluster is
uniformly separated from zero. Thus all $m$ cluster eigenvalues
are positive for $r<r_*$ close enough, and
\[
 0<A_r|_{\ran P_r}\le\varepsilon_r\Id,
 \qquad \varepsilon_r\longrightarrow0.
\]
The inverse on $\ran(\Id-P_r)$ is uniformly bounded.
This argument does not require $w_r$ to be pointwise monotone.

Choose an orthonormal basis $\phi_1,\ldots,\phi_m$ of
$\ker A_{r_*}$ and put $g_j=\partial_\nu\phi_j|_\T$.
Elliptic regularity gives $g_j\in C^\infty(\T)$.
They are linearly independent by uniqueness for zero Dirichlet and
Neumann data. In this setting that uniqueness follows by transferring
the equation to $\Omega_{r_*}$ and extending the solution by zero: both
vanishing traces make the extension solve the Helmholtz equation
distributionally on the plane, so its Fourier transform, an $L^2$
function supported on $\{|\xi|=k\}$, must vanish.

Set $h_j=g_j$ and choose fixed smooth extensions $\mathcal E h_j$ to
$\overline\D$. Extend $\mathcal E$ linearly on their span. For $h=h_c$ set
$f_{r,h}=\mathscr A_r\mathcal E h\in L^2(\D)$,
$v_{r,h}=A_r^{-1}f_{r,h}$, and $u_{r,h}=\mathcal E h-v_{r,h}$.
Then $u_{r,h}$ solves $\mathscr A_ru_{r,h}=0$ with boundary value $h$.
If
\[
 q_r[p,q]=\int_\D
 \bigl(\nabla\overline p\cdot\nabla q-k^2w_r\overline p\,q\bigr),
\]
Green's identity gives $q_r[v_{r,h},u_{r,h}]=0$ and
$\ip h{\Lambda_rh}=q_r[u_{r,h},u_{r,h}]$. Since
$\mathcal E h=u_{r,h}+v_{r,h}$ and
$q_r[v_{r,h},v_{r,h}]=\ip{f_{r,h}}{A_r^{-1}f_{r,h}}$, we obtain
\begin{equation}\label{eq:schur}
 \ip h{\Lambda_rh}
 =\int_\D\bigl(|\nabla\mathcal E h|^2-k^2w_r|\mathcal E h|^2\bigr)
       -\ip{f_{r,h}}{A_r^{-1}f_{r,h}}.
\end{equation}
At $r_*$, a second application of Green's identity yields
\[
 \ip{\phi_i}{f_{r_*,h}}=\int_\T\overline{g_i}h\,d\theta.
\]
For $h=\sum c_jg_j$, these coefficients are the product of the
positive-definite Gram matrix of the $g_j$ with $c$.
Continuity gives $\norm{P_rf_{r,h_c}}\ge c_0|c|$ for a fixed
$c_0>0$ and nearby $r$.
The cluster contribution to the last term of \eqref{eq:schur} is
at most $-c_0^2|c|^2/\varepsilon_r$. The remaining terms are
bounded in absolute value by $C|c|^2$. Taking
$M_r=c_0^2/\varepsilon_r-C$ proves \eqref{eq:pole}.
\end{proof}

We now use the boundary identity to transfer these negative
directions to the spectrum of the Cayley transform.

\begin{proposition}\label{prop:near}
Under the assumptions of Lemma~\ref{lem:pole}, for every
$\pi<\alpha<2\pi$ and every sufficiently close nonresonant
$r<r_*$, the operator $U_r$ has at least $m$ eigenvalues with
arguments in $(\alpha,2\pi)$, counted with multiplicity.
\end{proposition}

\begin{proof}
Let $J_0:\mathbb C^m\to H^{1/2}(\T)$ be the map $J_0c=h_c$.
At each fixed nonresonant $r$, Lemma~\ref{lem:density} provides
Herglotz traces $\widetilde h_{r,j}$ such that the map
$J_rc=\sum c_j\widetilde h_{r,j}$ satisfies
\[
 \norm{J_r-J_0}\le\delta_r\le1,\qquad
 \norm{\Lambda_r}\delta_r(2\norm{J_0}+1)\le1.
\]
Here $\norm{\Lambda_r}$ is the norm from $H^{1/2}$ to $H^{-1/2}$.
The quadratic-form error is at most $|c|^2$, so \eqref{eq:pole}
gives
\[
 \ip{J_rc}{\Lambda_rJ_rc}\le-(M_r-1)|c|^2,
 \qquad \norm{J_rc}_{H^{1/2}}\le C_0|c|,
\]
with $C_0$ independent of $r$. The approximation is chosen separately
at each radius; no uniform bound on the Herglotz densities is needed.

Set $x_{r,c}=\OO_r^*J_rc$. Proposition~\ref{prop:identity} places
these vectors in $\Dom\KK_r$. Equations \eqref{eq:identity},
\eqref{eq:remainder}, and \eqref{eq:Ostar} give
\[
 \ip{x_{r,c}}{\KK_rx_{r,c}}
       \le-\frac4{k^2}(M_r-C_1)|c|^2,
 \qquad \norm{x_{r,c}}\le C_2|c|,
\]
with constants uniform near $r_*$. When $M_r>C_1$, this map is
injective, since a zero image would contradict the strict negative
bound. Its image is an $m$-dimensional subspace of $\Dom\KK_r$
on which the Rayleigh quotient is at most
$-4(M_r-C_1)/(k^2C_2^2)$.

For any $D>0$, the spectral projection of $\KK_r$ onto
$(-\infty,-D)$ therefore has rank at least $m$ when $r$ is close
enough to $r_*$. Otherwise that image subspace would contain a
nonzero vector orthogonal to the projection, with quadratic form
at least $-D$ times its squared norm, a contradiction.
Choose $D$ so that $-D<\cot(\alpha/2)$, and use \eqref{eq:cot}.
\end{proof}

Only Herglotz waves have been differentiated to form infinite
sequences. No such sequence has been assigned to an arbitrary
Dirichlet eigenfunction, and no eigenvalue exactly equal to $1$
is asserted at a resonant radius.

\section{Counting and completion of the proof}\label{sec:completion}

\subsection{Domains with analytic boundary}
We retain the spectral tail in the counting argument. The energy
parameter in $L_{r,k}$ is written explicitly from this point onward.

\Needspace{12\baselineskip}
\begin{proposition}\label{prop:remainder}
Let $\Omega=F(\D)$, where $F$ is holomorphic and injective on a
neighborhood of $\overline\D$. For every $k>0$ and every integer
$m$ with $1\le m\le N_\Omega(k^2)$,
\begin{equation}\label{eq:count-fan}
 2\pi m\le\int_0^1\Phi_m(L_{r,k})\,dr,
\end{equation}
and
\begin{equation}\label{eq:quantitative-analytic}
 k^2|\Omega|\ge4\pi m+2\int_0^1\tau_m(L_{r,k})\,dr>4\pi m.
\end{equation}
\end{proposition}

\begin{proof}
Continuity and strict monotonicity of $\lambda_j(\Omega_r)$ were
proved in Lemma~\ref{lem:pole}. Since $w_r\le Cr^2$ near zero,
the disk Poincar\'e inequality gives
$\lambda_1(\Omega_r)\to\infty$ as $r\downarrow0$.
It follows that there are finitely many radii
$0<r_1<\cdots<r_q\le1$ where $k^2$ is a Dirichlet eigenvalue.
If their multiplicities are $m_1,\ldots,m_q$, then
\begin{equation}\label{eq:count}
 \sum_{\ell=1}^q m_\ell=N_\Omega(k^2).
\end{equation}
Indeed, each ordered eigenvalue at most $k^2$ at $r=1$ crosses
that value exactly once, and every other eigenvalue stays above it.
Under the assumptions of the proposition, $q\ge1$.

Choose nonresonant radii $a_1\in(0,r_1)$ and
$a_\ell\in(r_{\ell-1},r_\ell)$ for $\ell\ge2$.
Lemma~\ref{lem:phase}(ii) gives a common
$\alpha_0\in(\pi,2\pi)$ such that no $U_{a_\ell,k}$ has spectrum
with argument in $[\alpha_0,2\pi)$.
Fix $\alpha\in(\alpha_0,2\pi)$. Proposition~\ref{prop:near}
provides $b_\ell\in(a_\ell,r_\ell)$ at which there are at least
$m_\ell$ eigenvalues with arguments in $(\alpha,2\pi)$.
The selected eigenvalues persist in that open arc under small
perturbations of $b_\ell$. Since crossings of $e^{i\alpha}$ are
isolated, we may also arrange that
$e^{i\alpha}\notin\sigma(U_{b_\ell,k})$.

Lemma~\ref{lem:phase}(iii) gives at least $m_\ell$ crossings in
each $(a_\ell,b_\ell)$. These intervals are disjoint. By
\eqref{eq:count}, the full path has at least $N_\Omega(k^2)$
crossings, and hence at least $m$. Applying the strengthened
estimate \eqref{eq:fan-cost} to this prescribed number gives
\[
 \alpha m\le\int_0^1\Phi_m(L_{r,k})\,dr.
\]
The right-hand side is independent of $\alpha$. Letting
$\alpha\uparrow2\pi$ proves \eqref{eq:count-fan}.
Equations \eqref{eq:fan-def} and \eqref{eq:area} give
\eqref{eq:quantitative-analytic}; its last inequality follows from
\eqref{eq:positive-tail} and continuity. If $r_q=1$, the argument
still uses $b_q<1$, so the endpoint multiplicity is included.
\end{proof}

\subsection{Arbitrary boundary}
A limit of strict inequalities need not be strict. We therefore pass
the positive spectral tail, rather than only the strict inequality,
through the conformal exhaustion. This completes the proof of
Theorem~\ref{thm:main}.

\begin{proof}
Let $F:\D\to\Omega$ be a Riemann map and set
$\Omega_s=F(s\D)$ for $0<s<1$. The map $z\mapsto F(sz)$ is
holomorphic and injective on a neighborhood of $\overline\D$.
The boundary transport for $F$ is defined on each compact radial
interval in $[0,1)$; the constructions on these intervals agree by
uniqueness of the differential equations. We denote its generator
by $L_{r,k}$ for $0\le r<1$ and $k>0$.

Fix $j\ge1$. We first recall spectral convergence under this
exhaustion. Domain monotonicity shows that $\lambda_j(\Omega_s)$
decreases to a limit not smaller than $\lambda_j(\Omega)$.
Conversely, the min--max principle and the form-core property of
$C_c^\infty(\Omega)$ give, for every $\varepsilon>0$, a
$j$-dimensional subspace $V\subset C_c^\infty(\Omega)$ such that
\[
 \max_{0\ne u\in V}
 \frac{\int_\Omega|\nabla u|^2}{\int_\Omega|u|^2}
 \le\lambda_j(\Omega)+\varepsilon.
\]
The union of the supports of a basis of $V$ is compact in $\Omega$,
so it is contained in $\Omega_s$ when $s$ is sufficiently close to
$1$. The min--max principle on $\Omega_s$ then gives
$\lambda_j(\Omega_s)\le\lambda_j(\Omega)+\varepsilon$.
Thus, with
\begin{equation}\label{eq:exhaustion-energy}
 k_s=\sqrt{\lambda_j(\Omega_s)},\qquad
 k=\sqrt{\lambda_j(\Omega)},
\end{equation}
we have $k_s\to k>0$ and $|\Omega_s|\uparrow|\Omega|$.

Apply Proposition~\ref{prop:remainder} to $F(s\,\cdot)$ at energy
$k_s^2$ and with $m=j$. This is permitted since
$N_{\Omega_s}(k_s^2)\ge j$, including any multiplicity at that
energy. Under the radial change of variable $r=st$, the unitary
path for $F(s\,\cdot)$ is $U_{st,k_s}$ and its generator is
$sL_{st,k_s}$. This follows also for the initial-point correction
from the defining differential equation for $R_r$.
Since $\tau_j$ is positively homogeneous, the proposition gives
\begin{equation}\label{eq:exhaustion-remainder}
 |\Omega_s|\lambda_j(\Omega_s)
 \ge4\pi j+2\int_0^s\tau_j(L_{r,k_s})\,dr.
\end{equation}

For every fixed $r\in(0,1)$, the joint trace-norm continuity in
Proposition~\ref{prop:area} and \eqref{eq:fan-continuity} imply
\[
 \tau_j(L_{r,k_s})\longrightarrow\tau_j(L_{r,k})
 \qquad(s\uparrow1).
\]
Choose any sequence $s_n\uparrow1$. Apply Fatou's lemma on $(0,1)$
to the nonnegative functions
$\mathbf1_{(0,s_n)}(r)\tau_j(L_{r,k_{s_n}})$, and use
\eqref{eq:exhaustion-remainder} and \eqref{eq:exhaustion-energy}.
We obtain
\begin{equation}\label{eq:quantitative-rough}
 |\Omega|\lambda_j(\Omega)
 \ge4\pi j+2\int_0^1\tau_j(L_{r,k})\,dr.
\end{equation}
The integral is finite because
\[
 0\le\tau_j(L_{r,k})\le\tr L_{r,k},\qquad
 \int_0^1\tr L_{r,k}\,dr=\frac{k^2}{2}|\Omega|;
\]
the trace identity follows first on compact radial intervals and
then by monotone convergence. The integral in
\eqref{eq:quantitative-rough} is strictly positive. Indeed,
$L_{r,k}$ is injective on $\ell^2(\mathbb N_0)$ for every $r>0$,
and $\tau_j(L_{r,k})$ is positive and continuous on every compact
subinterval of $(0,1)$. Thus \eqref{eq:quantitative-rough} proves
\eqref{eq:strict-eigenvalues}.

Finally, let $E>0$. If $N_\Omega(E)=0$, then \eqref{eq:polya} is
immediate. Otherwise put $M=N_\Omega(E)$. Since
$\lambda_M(\Omega)\le E$, the strict eigenvalue inequality gives
\[
 4\pi M<|\Omega|\lambda_M(\Omega)\le|\Omega|E.
\]
This proves \eqref{eq:polya} and completes the proof.
\end{proof}

The remainder in \eqref{eq:quantitative-rough} depends on the domain,
the chosen conformal parametrization, and the index. Its positivity
is sufficient for the strict inequality; no uniform lower bound
in terms of the index alone is asserted here.

\section{Lean formalization}\label{sec:lean-formalization}

Both assertions of Theorem~\ref{thm:main} have been formalized in Lean
in the modular source tree rooted at \texttt{RequestProject/Main.lean}.
The aggregate file imports the three checked modules
\texttt{MainBase.lean}, \texttt{MainMiddle.lean}, and
\texttt{MainSpectral.lean}. The declarations
\texttt{main} and \texttt{main\_counting} prove
\eqref{eq:strict-eigenvalues} and \eqref{eq:polya}, respectively, for
every bounded nonempty simply connected planar domain, without a
boundary regularity assumption. The formalization identifies
$\mathbb C$ with $\mathbb R^2$ using the usual Euclidean norm and
two-dimensional Lebesgue measure.

The development constructs the real and complex spaces $H_0^1(\Omega)$
from the smooth compactly supported core, and the self-adjoint
Dirichlet operator associated with the closed gradient-energy form
$q_\Omega[u]=\int_\Omega |\nabla u|^2$. Its positive eigenvalues are
defined independently of the smooth-core min--max, using orthonormal
eigenfamilies of the compact inverse. Lean proves that the two
constructions agree, together with the nondecreasing ordering,
divergence to infinity, and correct real and complex multiplicities.
Thus the formalized counting function is precisely
$N_\Omega(E)=\#\{j\ge1:\lambda_j(\Omega)\le E\}$, with the endpoint
included and eigenvalues repeated according to multiplicity.

The source tree is pinned to
\texttt{leanprover/lean4:v4.35.0-rc2} and the matching \texttt{mathlib}
revision recorded in the repository's \texttt{lake-manifest.json}.
The Palomar preparation uses this pinned toolchain; the repository records
the exact build and audit commands.
using \texttt{lake build +RequestProject.Main}. Axiom audits of both
final declarations and the operator and spectral identification
theorems report only \texttt{propext}, \texttt{Classical.choice}, and
\texttt{Quot.sound}, the standard principles of Lean's classical
mathematics~\cite{LeanAxioms}. They do not depend on \texttt{sorryAx}
or additional axioms; the source contains no \texttt{sorry},
\texttt{admit}, or \texttt{native\_decide} proof steps.

The original variational development was produced with the assistance
of Aristotle; the operator and spectral identification was developed
with the assistance of Codex. The source, pinned dependencies,
verification commands, build record, and source checksum are provided
in the accompanying GitHub repository.\footnote{\url{https://github.com/QuanyuTang/polya-dirichlet-simply-connected-formalization-Palomar-version}.
The repository is currently private; access must be granted by its owner.}

\begin{thebibliography}{99}

\bibitem{Berezin}
F.~A. Berezin,
\emph{Covariant and contravariant symbols of operators},
Math. USSR-Izv. \textbf{6} (1972), 1117--1151.

\bibitem{CK}
D. Colton and R. Kress,
\emph{On the denseness of Herglotz wave functions and electromagnetic
Herglotz pairs in Sobolev spaces},
Math. Methods Appl. Sci. \textbf{24} (2001), 1289--1303.

\bibitem{EP}
J.-P. Eckmann and C.-A. Pillet,
\emph{Spectral duality for planar billiards},
Comm. Math. Phys. \textbf{170} (1995), 283--313.

\bibitem{Fan}
K. Fan,
\emph{On a theorem of Weyl concerning eigenvalues of linear
transformations. I},
Proc. Natl. Acad. Sci. USA \textbf{35} (1949), 652--655.

\bibitem{FLPSballs}
N. Filonov, M. Levitin, I. Polterovich, and D.~A. Sher,
\emph{P\'olya's conjecture for Euclidean balls},
Invent. Math. \textbf{234} (2023), 129--169.

\bibitem{FLPSannuli}
N. Filonov, M. Levitin, I. Polterovich, and D.~A. Sher,
\emph{P\'olya's conjecture for Dirichlet eigenvalues of annuli},
J. Lond. Math. Soc. (2) \textbf{113} (2026), e70425.

\bibitem{FS}
P. Freitas and I. Salavessa,
\emph{Families of non-tiling domains satisfying P\'olya's conjecture},
J. Math. Phys. \textbf{64} (2023), 121503.

\bibitem{Kato}
T. Kato,
\emph{Perturbation Theory for Linear Operators},
2nd ed., reprint, Classics in Mathematics,
Springer-Verlag, Berlin, 1995.

\bibitem{LiYau}
P. Li and S.-T. Yau,
\emph{On the Schr\"odinger equation and the eigenvalue problem},
Comm. Math. Phys. \textbf{88} (1983), 309--318.

\bibitem{LeanAxioms}
The Lean developers,
\emph{The Lean Language Reference}, Section~8, Axioms,
\url{https://lean-lang.org/doc/reference/latest/Axioms/}
(accessed 3 October 2026).

\bibitem{Polya54}
G. P\'olya,
\emph{Mathematics and Plausible Reasoning},
Vols.~I--II, Princeton University Press, Princeton, NJ, 1954.

\bibitem{Polya61}
G. P\'olya,
\emph{On the eigenvalues of vibrating membranes},
Proc. London Math. Soc. (3) \textbf{11} (1961), 419--433.

\bibitem{Weyl}
H. Weyl,
\emph{Das asymptotische Verteilungsgesetz der Eigenwerte linearer
partieller Differentialgleichungen (mit einer Anwendung auf die
Theorie der Hohlraumstrahlung)},
Math. Ann. \textbf{71} (1912), 441--479.

\end{thebibliography}
\end{document}
